\chapter*{Introduction générale}
\addcontentsline{toc}{chapter}{Introduction générale}

La transformation numérique des fonctions \textit{Talent Acquisition} ne se limite plus à la publication d'offres et au suivi administratif des candidatures. Dans les contextes de recrutement à volume élevé, les équipes doivent simultanément centraliser des CV hétérogènes, structurer les données utiles, coordonner plusieurs intervenants métiers, documenter chaque décision et accélérer les délais de traitement sans sacrifier la qualité d'évaluation. Cette pression opérationnelle crée un terrain favorable à l'automatisation, à l'analytique temps réel et à l'intelligence artificielle appliquée aux flux RH.

C'est dans ce cadre que s'inscrit \projectname{}, une plateforme de recrutement orientée entreprise développée avec \sourcecode{Next.js 16}, \sourcecode{TypeScript}, \sourcecode{Bun}, \sourcecode{Neon PostgreSQL}, \sourcecode{Drizzle ORM} et un socle d'IA s'appuyant sur des API compatibles OpenAI/NVIDIA. Les documentations officielles des principales technologies citées sont regroupées dans la netographie. La solution couvre l'ensemble du cycle de recrutement : page d'accueil et authentification, gestion d'un vivier de CV, création et suivi des postes, assignation des candidats, historique des transitions, planification des entretiens, rapports d'évaluation, notifications, exports, tableaux de bord administratifs, centre de gouvernance et assistant conversationnel outillé.

L'originalité majeure de la solution réside dans sa couche d'intelligence. Le système ne se contente pas d'ajouter une fonctionnalité d'aide à la rédaction ; il intègre un agent conversationnel multi-outils, un pipeline hybride de recherche \textit{Retrieval-Augmented Generation} (RAG), des mécanismes de \textit{grounding} destinés à limiter les réponses non fondées, une confirmation explicite avant les actions IA qui modifient les données et des évaluations quantitatives sur un jeu de 40 requêtes. Le principe RAG s'appuie sur la littérature de référence en récupération augmentée, tandis que l'implémentation du projet l'adapte aux contraintes métier RH.

L'objectif de ce rapport est de présenter, dans un cadre académique, la conception et la réalisation de cette plateforme selon une progression structurée : un cadrage général, un chapitre d'analyse et d'architecture, puis cinq chapitres d'implémentation organisés par modules fonctionnels. Cette structuration permet d'exposer à la fois la logique métier, les décisions techniques et la valeur produite pour les différents rôles de l'application : \textbf{TA}, \textbf{Manager}, \textbf{HR} et \textbf{Admin}.

La question directrice du mémoire peut donc être formulée ainsi : comment concevoir une plateforme de recrutement qui centralise les opérations métier tout en intégrant une assistance IA utile, mesurable et gouvernée ? Cette question guide la structure du rapport : chaque chapitre ne présente pas seulement une fonctionnalité isolée, mais une brique nécessaire à cette chaîne complète, depuis l'accès sécurisé jusqu'à la supervision des décisions et des communications.

Le présent mémoire est organisé comme suit :
\begin{itemize}
  \item le \textbf{chapitre 1} présente l'entreprise d'accueil, le contexte métier du recrutement et les objectifs du projet ;
  \item le \textbf{chapitre 2} formalise les besoins, le Product Backlog, les objectifs de sprint, l'architecture applicative, le modèle de données et l'environnement technique ;
  \item le \textbf{chapitre 3} traite le module d'accès, d'identité et d'entrée par rôle ;
  \item le \textbf{chapitre 4} couvre la gestion du vivier de CV et des offres d'emploi ;
  \item le \textbf{chapitre 5} détaille le pipeline candidat et la conduite des entretiens ;
  \item le \textbf{chapitre 6} expose la couche d'intelligence artificielle, le matching, l'autopilot d'entretien et le pipeline RAG ;
  \item le \textbf{chapitre 7} présente la supervision, les notifications, l'analytique et la gouvernance opérationnelle ;
  \item enfin, une \textbf{conclusion générale} synthétise les acquis, la validation et les apports du projet.
\end{itemize}

Ainsi, ce rapport ne décrit pas seulement une interface web de recrutement, mais une plateforme métier complète, conçue pour fiabiliser les décisions, améliorer la traçabilité des opérations et préparer l'industrialisation d'usages IA dans un contexte RH exigeant.
